% Part: history
% Chapter: biographies
% Section: bertrand-russell

\documentclass[../../../include/open-logic-section]{subfiles}

\begin{document}

\olfileid{his}{bio}{rus} 

\olsection{برټرېنډ رسل}

\olphoto{russell-bertrand}{برټرېنډ رسل}

برټرېنډ رسل د نوې تجزيه‌يي فلسفې د بنسټ
ايښودونکو له ډلې ګڼل کېږي۔ رسل د 1872
د مې پر 18مه زېږېدلے و۔ هغه يوازې په
فلسفه او منطق کښې د خپل کار له امله
نۀ پېژندل کېده، بلکې په بېلابېلو
موضوعاتو ئې ډېر عام‌فهم کتابونه هم
ليکلي وو۔ هغه د خپل ټول ژوند په
اوږدو کښې يو تود سياسي فعال هم و۔

رسل د وېلز د مونماوتشاير په ټرېلېک کښې
زېږېدلے و۔ مور او پلار ئې د برتانوي
اشرافو غړي وو۔ هغوي آزادخياله وو،
او د هغه وخت د بوسټن له تندلارو سره
ئې هم ملګرتيا کړې وه۔ له بده مرغه،
د رسل مور او پلار د هغه په کم عمرۍ
کښې وفات شول، او رسل د خپل نيکۀ
او نيا سره اوسېدو ته ولېږل شو۔ هلته
ئې مذهبي روزنه وشوه؛ دا هغه څه وو
چې مور او پلار ئې په هر قيمت ترې
ځان ساتل غوښتل۔ نيا ئې د اخلاقو
په ټولو چارو کښې ډېره سخته وه۔ په
تنکۍ ځوانۍ کښې ئې زياتره په کور
کښې له شخصي ښوونکو زده کړه کوله۔

په تجزيه‌يي فلسفه، او په تېره په منطق،
کښې د رسل اغېز ډېر ستر دے۔ هغه د
کېمبرېج په ټرينيټي کالج کښې رياضيات
او فلسفه ولوستل، او هلته د رياضيپوه
او فيلسوف الفرېډ نورټ وېټ هېډ تر
اغېز لاندې راغے۔ په 1910 کښې رسل
او وېټ هېډ د \emph{د رياضياتو اصول}
لومړے ټوک خپور کړ، او پکښې ئې د
دې نظر ملاتړ وکړ چې رياضيات منطق
ته راکمول کېدلے شي۔ وروسته ئې په
سلګونو کتابونه، مقالې او سياسي
رسالې خپرې کړې۔ په 1950 کښې ئې
د ادبياتو نوبېل جايزه وګټله۔

رسل په سياست او ټولنيز فعاليت کښې ژور
ښکېل و۔ د لومړۍ نړيوالې جګړې پر
وخت، د سوله‌پالو فعاليتونو او احتجاج
له امله ونيول شو او د شپږو مياشتو
دپاره زندان ته ولېږل شو۔ په زندان
کښې ئې ليکل او لوستل کولے شول،
او وايي چې دا تجربه ئې «خورا د
خوښې وړ» وموندله۔ هغه د خپل ټول
ژوند په اوږدو کښې سوله‌پال پاتې شو،
او په 1961 کښې د اټومي وسلو د
لرې کولو په يوه غونډه کښې د ګډون
له امله بيا بندي شو۔ هغه په 1948
کښې د الوتکې له يوې پېښې هم ژوندے
پاتې شو؛ په دې پېښه کښې يوازې د
سګرټ څښونکو په برخه کښې ناست کسان
ژوندي پاتې شوي وو۔ له دې امله رسل
ويل چې ژوند ئې د سګرټ څښلو مرهون
دے۔ رسل څلور ځله وادۀ کړے و، خو
د وادۀ څخه بهر د عاشقانه اړيکو
شهرت ئې هم درلود۔ هغه د 1970
د فبرورۍ پر 2مه د 97 کلونو په
عمر د وېلز په پېنريندېودراېت کښې
وفات شو۔ \emph{سرچينه‌يي يادونه
OLSTH-044:} د ټول ژوند د سوله‌پالۍ
پورته مطلقه دعوه د منجمد متن ده۔
د مکماسټر د رسل آرشيف تاريخچه او
د هغه د راټولو ليکنو فهرست ښيي
چې هغه وروسته د هټلر پر ضد د
دويمې نړيوالې جګړې ملاتړ وکړ۔ نو
د سولې د ملاتړ اوږد فعاليت ئې
د هرې جګړې له بې‌استثنا ردولو
سره يو شان نۀ دے۔

\begin{reading}
رسل په درېو برخو کښې خپل ژوندليک
وليکۀ، چې له 1872--1967 پورې ئې
ژوند رانغاړي
\citep{Russell1967,Russell1968,Russell1969}۔
په مکماسټر پوهنتون کښې د برټرېنډ
رسل د څېړنې مرکز د هغه د آرشيف
کور دے۔ د راټولو آثارو د ټوکونو
په هکله د معلوماتو، د لټون وړ
فهرستونو او آرشيفي پروژو دپاره
ئې په \citet{Duncan2015} کښې وېبپاڼه
وګورئ۔ د رسل مقاله \emph{د دلالت
په هکله} \citep{Russell1905} د
شلمې پېړۍ د تجزيه‌يي فلسفې له
کلاسيکو آثارو ده۔

د فلسفې د سټېنفورډ پوهنغونډ په
\citep{Irvine2015} کښې د رسل مدخل
د خواهش او سياسي تيورۍ په هکله
د هغه د خبرو غږيزې ټوټې لري۔ د
رسل ډېرې ويډيويي مرکې په انټرنېټ
کښې شته۔ د بېلګې په توګه، د
سګرټ څښلو او د الوتکې په پېښه
کښې د ښکېلتيا په هکله ئې خبرې
په \citet{RussellND} کښې وګورئ۔
د هغه ځينې آثار، چې \emph{د
رياضيکي فلسفې پېژندنه} هم پکښې
ده، په \citet{LibriVoxND} کښې د
وړيا غږيزو کتابونو په توګه شته۔
\end{reading}

\end{document}
